\documentclass[12pt]{article}
\usepackage[utf8]{inputenc}
\usepackage[T1]{fontenc}
\usepackage{hyperref}
\title{Wiru Language Resources for K-12 Curriculum Development}
\author{Compiled by Claude Sonnet 5.5 (claude-sonnet-5-5) for Andr\'as Kornai}
\date{}
\begin{document}
\maketitle

\section{Introduction}
Wiru (ISO 639-3 wiu) is spoken in the Southern Highlands of Papua New Guinea, in the Ialibu--Pangia area. Early counts gave about 15,000 speakers, and a recent ANU project estimates 15,000--40,000. Sources disagree on classification. Wikipedia places it in Teberan--Pawaian, and Joshua Project describes it as Trans-New Guinea. My searches found very little material that is usable in schools. Wiru is not to be confused with Witu or Vitu (wiv), which is a different language.

\section{Primary Educational Resources}
I found no Wiru primers, textbooks or school curricula in my searches. Tok Pisin is reported as the language of school and trade.

\section{Linguistic and Reference Materials}
Harland Kerr's 1967 grammar (a University of Hawai'i MA thesis), his 1972 dictionary and his 1975 comparative study are the main scholarly sources. Search results mentioned them, but I did not see a URL for any of them, so they are not in the JSON. General description: \url{https://en.wikipedia.org/wiki/Wiru_language}.

An ANU research project on children's acquisition of Wiru describes the language as under-documented: \url{https://chl.anu.edu.au/event/childrens-acquisition-wiru-trans-new-guinea-language-thesis-proposal-review}.

\section{Religious and Community Resources}
Joshua Project reports that Wiru scripture portions and audio Bible resources exist, and it gives a community profile: \url{https://joshuaproject.net/people_groups/15402}. I did not see direct links to the scripture texts themselves.

\section{Digital and Audio Resources}
I found no dedicated Wiru digital or audio language resources beyond the audio Bible mentioned above.

\section{Curriculum Development Context}
Little is publicly available for Wiru. A curriculum would have to start from Kerr's grammar and dictionary, scripture portions, and new community-based work with Wiru speakers, with Tok Pisin and English as bridge languages.
\end{document}